\documentclass[11pt]{article}
\usepackage[margin=1in]{geometry}
\usepackage{amsmath,amssymb,amsthm}
\newtheorem{theorem}{Theorem}
\newtheorem{lemma}{Lemma}
\newtheorem{corollary}{Corollary}
\newcommand{\Prob}{\mathbb P}
\newcommand{\E}{\mathbb E}
\newcommand{\supp}{\operatorname{supp}}
\title{Small-support exclusion and exact near-point closure\\for mixed bounded degrees}
\author{Independent support reconstruction; no historical novelty claim}
\date{November 15, 2026}
\begin{document}
\maketitle

\paragraph{Scope.}
This note proves only the small-support and near-point parts of
bounded-frame recovery. It allows unequal total degrees, no internal
edges, and internal degree one. It needs neither a symmetric omission
weight nor the strict-weight hypothesis used by a separate entropy
argument. It does not assert macroscopic rigidity, a polynomial
intertwiner moment, or recovery for every nonregular frame.

\section{Typed graph model and uniform constants}
Fix $K\ge1$ and $D\ge3$. There are $1\le k\le K$ planted blocks
$V_i$, with positive sizes $n_i$ and total size $n=\sum_i n_i$.
The nonnegative integer frame $Q=(d_{ij})$ has total degrees
\[
 d_i=\sum_jd_{ij},\qquad 3\le d_{\min}:=\min_i d_i\le d_i\le D.
\]
The off-diagonal support is connected, with the single-type case
allowed. Undirected balance and admissibility mean
\[
 n_i d_{ij}=n_jd_{ji},\qquad n_i d_{ii}\ \text{even}.
\]
An absent factor is omitted. An internal factor pairs its $n_id_{ii}$
labelled stubs uniformly. A cross factor independently matches its
two equal pools of $n_id_{ij}=n_jd_{ji}$ labelled stubs uniformly.
Until simplicity conditioning, loops and parallel edges are allowed;
a loop contributes two to the adjacency diagonal and degree.

Put
\[
 \rho=\frac1{K D^{K-1}}.
\]
Every block has $n_i\ge\rho n$. Indeed, along a frame path each
positive ratio of block sizes lies between $D^{-1}$ and $D$, and a
path has length at most $k-1$. Every positive stub pool therefore
has at least $\rho n$ stubs. Connectedness fixes the size ratios;
with fixed $K,D$ there are finitely many possible frames and ratios.

Call a nonempty set $S$ \emph{bad} if no vertex outside $S$ has
exactly one incident edge into $S$, counting edge multiplicities.

\begin{theorem}[A single simultaneous graph event]\label{thm:small}
There is an explicit $\eta=\eta(K,D)>0$ such that, in the model with
all factors conditioned simple,
\[
 \Prob\bigl(\exists\text{ bad }S:\ 1\le|S|\le\eta n\bigr)
 =O_{K,D}\bigl(n^{-(d_{\min}-2)}\bigr).
\]
On the complementary event, for every field $F$, every integer
$r\ge1$, every $R\in F^{r\times r}$, and every $V\in F^{n\times r}$,
\[
 AV=VR,\qquad |\{v:V_v\ne0\}|\le\eta n
 \quad\Longrightarrow\quad V=0.
\]
All these fields, dimensions and matrices are handled simultaneously
by this one graph event; there is no union bound over coefficients.
\end{theorem}

\section{A full typed witness count}
First work in the unconditioned stub model. Let $S$ be bad, and write
\[
 s=|S|,\quad W=N(S)\setminus S,\quad w=|W|.
\]
Let $a$ be the number of edges internal to $S$, counting a loop once,
and $b$ the number of edges from $S$ to $W$. Put $e=a+b$ and
$H_S=\sum_{v\in S}d_v$. Then
\[
 2a+b=H_S\ge d_{\min}s,\qquad b\ge2w,
 \qquad e\le Ds,\quad w\le Ds/2.
\]
In particular, for $c=(d_{\min}-2)/2\ge1/2$,
\begin{equation}\label{eq:excess}
 e-s-w=\frac{H_S}2-s+\frac b2-w\ge cs,
 \qquad e-s-w\ge\lceil cs\rceil.
\end{equation}
This is the only place where the lower bound on the \emph{varying}
total degrees enters the counting exponent.

A witness specifies $S,W$, every stub at $S$, the subset of stubs at
$W$ paired to $S$, and all $e$ resulting pairs. Every actual bad set
has such a witness. For $s\le\rho n/(4D)$ and an admissible witness,
a prescription of $m$ pairs in an internal pool of size $L$ has
probability
\[
 \frac1{(L-1)(L-3)\cdots(L-2m+1)}.
\]
In a cross pool of size $L$ on each side the probability is
$1/(L)_m$. There are at most $e\le Ds$ prescribed pairs in any
factor; every denominator factor is consequently at least $\rho n/2$.
By independence across factors, the full witness probability is at most
\begin{equation}\label{eq:prob}
 (2/(\rho n))^e.
\end{equation}
Prescriptions that violate types or pair a stub twice have probability
zero and may be included in an upper count.

For fixed $s,w,e$, an upper bound for the number of witnesses is
\begin{equation}\label{eq:witness}
 \binom ns\binom nw\,2^{Dw}(2Ds)^e.
\end{equation}
Indeed the $W$ vertices have at most $Dw$ stubs to select from. The
selected stubs together with all stubs at $S$ number exactly $2e$ for
an actual witness, and their pairings number $(2e-1)!!\le(2e)^e
\le(2Ds)^e$. If a selected stub set has a different size, count no
pairings for that set. Formula~\eqref{eq:witness} also overcounts
choices with $S\cap W\ne\varnothing$ and pairs not touching $S$;
neither overcount invalidates the upper bound.

Combining \eqref{eq:prob}--\eqref{eq:witness}, and using
$\binom ns\le(\mathrm e n/s)^s$, gives at most
\[
 \exp(s+w)\,2^{Dw}(4D/\rho)^e
       (s/n)^{e-s-w}(s/w)^w
\]
in expected witnesses. Interpret $(s/w)^w$ as one for $w=0$.
For $w>0$, $w\log(s/w)\le s/\mathrm e$. Using
\eqref{eq:excess}, $w\le Ds/2$ and $e\le Ds$, and summing over
at most $(Ds+1)^2$ choices of $e,w$, proves
\begin{equation}\label{eq:Zs}
 \Prob(\exists\text{ bad }S:|S|=s)
 \le \E Z_s\le C^s(s/n)^{\lceil cs\rceil},
\end{equation}
where $Z_s$ counts all witnesses and one explicit choice is
\[
 C=4(D+1)^2
 \exp(1+D/2+1/\mathrm e)\,2^{D^2/2}(4D/\rho)^D.
\]
The rough inequality $(Ds+1)^2\le[4(D+1)^2]^s$, $s\ge1$, justifies
the displayed constant. The count permits every support composition
across the planted blocks.

Choose, uniformly over all values of $d_{\min}$,
\begin{equation}\label{eq:eta}
 0<\eta\le\min\left\{
    \frac\rho{4D},\frac1{4\mathrm e C^2},
    \frac\rho{2(D+1)}\right\}.
\end{equation}
For $u_s=C^s(s/n)^{cs}$, whenever $s+1\le\eta n$,
\[
 \frac{u_{s+1}}{u_s}
 \le C\mathrm e^c\eta^c=C(\mathrm e\eta)^c
 \le C(1/(4C^2))^{1/2}=\frac12.
\]
Here $c\ge1/2$ and $1/(4C^2)<1$; this explains why the same $\eta$
works for all mixed degree sequences. Since $(s/n)^{\lceil cs\rceil}
\le(s/n)^{cs}$, summing \eqref{eq:Zs} for $s\ge2$ gives
\begin{equation}\label{eq:config-tail}
 \Prob(\exists\text{ bad }S:2\le|S|\le\eta n)
 \le2C^2(2/n)^{d_{\min}-2}.
\end{equation}
An empty sum is understood as zero.

\section{Simplicity conditioning, including zero and one degrees}
For an internal degree $a_i=d_{ii}=0$ the factor is absent. For
$a_i=1$ it is a perfect matching and is simple. For $a_i\ge2$, the
expected numbers of loops and unordered double-edge pairs tend to
$(a_i-1)/2$ and $(a_i-1)^2/4$. These formulas follow respectively
from selecting two stubs at one vertex, or two stubs at each of two
vertices with either of the two cross pairings, and dividing by the
appropriate full-pairing denominator. For a cross $(u,v)$ factor,
the double-edge mean tends to $(u-1)(v-1)/2$. In particular, if
either degree is one the cross factor is automatically simple.

Here is a direct justification that these means control the zero-event.
For each fixed factorial-moment order, vertex-disjoint elementary bad
patterns give the product of the stated means in the limit. A
consistent overlapping collection of distinct patterns has at least one
extra edge over its number of vertices. This includes three parallel
edges formed by overlapping double-edge pairs, or two loops at a
single vertex. For a union on $v$ vertices with $e$ prescribed edges,
there are $O_{K,D}(n^v)$ choices of vertices and $O_{K,D}(n^{-e})$
prescription probability. Such overlaps therefore contribute $O(1/n)$.
Inconsistent stub prescriptions contribute zero. Fixed degrees and
$n_i\ge\rho n$ make all these bounds uniform. Finite Bonferroni
upper and lower sums for no elementary bad pattern, followed by an
increasing fixed truncation order, give the Poisson zero-event limit.

By independence of factors, their joint simplicity probability tends to
$\exp(-\Lambda)$, where
\[
 \Lambda=
 \sum_{i:d_{ii}\ge1}\frac{d_{ii}^2-1}4
 +\sum_{i<j:d_{ij}>0}\frac{(d_{ij}-1)(d_{ji}-1)}2
 \le K^2D^2.
\]
Internal degree zero contributes zero, rather than the negative value
that substitution into the first formula would produce. Thus there is
an eventual uniform lower bound
\[
 p_0=\tfrac12\exp(-K^2D^2)>0.
\]
There are finitely many admissible connected frames under fixed $K,D$,
each fixing positive size ratios, so the eventual size threshold can
also be uniform over this finite family.

On simplicity, a singleton is never bad: it has at least three distinct
neighbours, each with exactly one edge into it. Conditioning
\eqref{eq:config-tail} on all factors being simple therefore gives
\[
 \Prob_{\rm simple}(\exists\text{ bad }S:1\le|S|\le\eta n)
 \le p_0^{-1}2C^2(2/n)^{d_{\min}-2}.
\]
Finally each simple internal factor has the same $(d_{ii}!)^{n_i}$
labelled-stub multiplicity, and each simple cross factor has the same
$(d_{ij}!)^{n_i}(d_{ji}!)^{n_j}$ multiplicity. Thus conditioning gives
precisely the intended independent uniform simple factors, rather than
an unspecified contiguous distribution.

\section{Simultaneous algebraic and near-point closure}
\begin{proof}[Completion of Theorem~\ref{thm:small}]
Let the graph have no bad nonempty set of size at most $\eta n$, and
suppose $AV=VR$ over any field. If the nonzero row support $S$ is
nonempty and small, there is a vertex $v\notin S$ with exactly one
edge to a vertex $u\in S$. Then
\[
 (AV)_v=V_u\ne0,\qquad (VR)_v=V_vR=0,
\]
a contradiction. The coefficient is exactly one, so this works in
every characteristic and for every number of columns. The graph
event itself does not depend on any choice of field or candidate.
\end{proof}

Let $J$ be the planted $n$-by-$k$ indicator, so $AJ=JQ$ as an
integer identity. Let $Y$ be any one-hot $n$-by-$r$ matrix and
$R$ any integer $r$-by-$r$ matrix for which the \emph{integer}
identity $AY=YR$ has been verified. Let $T$ be any one-hot
$k$-by-$r$ matrix; the associated map may merge old types and need
not use all candidate labels.

\begin{corollary}[Exact closure near every point map simultaneously]
On the graph event in Theorem~\ref{thm:small},
\[
 |\{v:Y_v\ne(JT)_v\}|\le\eta n
 \quad\Longrightarrow\quad QT=TR\ \text{over the integers},
 \qquad Y=JT.
\]
No lower bound on the size of any candidate class is required.
\end{corollary}
\begin{proof}
Let $S$ be the changed row set. Its closed neighbourhood has at most
$(D+1)|S|\le\rho n/2<n_i$ vertices, by \eqref{eq:eta}. In every
old class $V_i$ choose $v_i$ outside that closed neighbourhood.
Its row and all neighbour rows are unchanged, giving
\[
 (QT)_{i,*}=(AJT)_{v_i,*}=(AY)_{v_i,*}
           =(YR)_{v_i,*}=(TR)_{i,*}.
\]
This proves integer compatibility. Subtracting gives
\[
 V:=Y-JT,\qquad AV=YR-JQT=YR-JTR=VR.
\]
Its row support is exactly $S$. Apply Theorem~\ref{thm:small} over
$\mathbb Q$ (or $\mathbb R$) to obtain $V=0$. If $S$ is empty,
the same clean-vertex argument also proves compatibility directly.
\end{proof}

\paragraph{Consequence for a separate recovery proof.}
If a macroscopic entropy theorem excludes all verified integer one-hot
candidates at positive distance from every point map, this corollary
closes the remaining candidates, including candidates with tiny new
classes or unused labels. The present graph event and closure hold
without equal total degrees and without assumptions on the candidate
matrix's eigenvalues, row injectivity, or nonsingularity. An exact
enumeration algorithm still needs its own unconditional modular-moment
bound; that bound is not supplied by small-support exclusion.
\end{document}
